\subsection{Results}

Mood fit averaged 4.71 under Vibe and 4.15 under Random
(Table~\ref{tab:pilot}). The paired difference was $+0.56$ rating points, with
a 95\% confidence interval from $-0.21$ to $1.33$ and a standardized paired
effect of $d_z=0.40$. Eight participant/session means were higher under Vibe,
one was tied, and six were lower. The retained directional paired test yielded
$t(14)=1.56$, $p=.070$; the supplied signed-rank sensitivity check yielded
$W^+=74.0$, $p=.088$. The interval includes effects on both sides of zero.

Liking averaged 4.56 under Vibe and 4.47 under Random. The paired difference
was $+0.09$, with a 95\% confidence interval from $-0.52$ to $0.71$ and
$d_z=0.08$. The directional tests yielded $t(14)=0.32$, $p=.375$, and
$W^+=58.5$, $p=.353$. No equivalence margin was specified, so the near-zero
estimate does not establish equivalence.

\begin{table*}[t]
  \caption{Exploratory summaries of 15 participant/session five-trial means.
  Confidence intervals are two-sided; retained tests use one-sided
  Vibe-over-Random alternatives.}
  \label{tab:pilot}
  \centering
  \small
  \begin{tabular}{@{}lcccccc@{}}
    \toprule
    Outcome & Vibe $M$ (SD) & Random $M$ (SD) &
    Difference [95\% CI] & $t(14)$, $p$ & $W^+$, $p$ & $d_z$ \\
    \midrule
    Mood fit & 4.71 (0.99) & 4.15 (1.30) &
    $+0.56$ [$-0.21,1.33$] & 1.56, .070 & 74.0, .088 & 0.40 \\
    Liking & 4.56 (1.18) & 4.47 (0.77) &
    $+0.09$ [$-0.52,0.71$] & 0.32, .375 & 58.5, .353 & 0.08 \\
    \bottomrule
  \end{tabular}
\end{table*}

\begin{figure}[t]
  \centering
  \includegraphics[width=\columnwidth]{fig3_pilot_evidence.pdf}
  \caption{\textbf{Paired participant/session outcomes.} Panels (a)--(b) show
  mood fit; panels (c)--(d) show liking. Left panels connect Random and Vibe
  five-trial means and summarize condition means. Right panels show all paired
  differences, fixed-seed bootstrap mean distributions, and reported paired
  means with $t$-based 95\% CIs. Both paired intervals cross zero.}
  \Description{Four compact panels compare mood fit and liking for 15
  participant-session means. Connected pairs show Random and Vibe ratings;
  estimation panels show paired differences and mean confidence intervals.}
  \label{fig:pilot}
\end{figure}

The pilot does not answer RQ2 decisively. Mood-fit ratings were higher on
average under Vibe Shuffle, whereas liking changed little, but neither estimate
excludes zero or establishes a recommendation benefit.

The paired display also shows heterogeneity that a condition mean would hide.
Several participant/session differences exceed one rating point in either
direction. That variation could reflect music familiarity, starting context,
order, signal availability, or genuine differences in how useful matching is.
The retained means cannot separate these explanations, but they argue against
treating one average response as a universal personalization effect.

\subsection{Reflection and design refinement}

The study and subsequent artifact audit clarify what the project can support.
Personal references provide a local comparison point, but they do not validate
the facial or cardiac mappings. A trajectory makes the decision easier to
inspect, but its arithmetic mean gives more weight to periods with more
coordinate changes. The matched Random condition controls catalog size,
playback path, exposure count, and rating procedure; selected songs can still
differ in artist, genre, language, familiarity, and popularity.

The audit produced concrete refinements to the research specification.
Emotion-detection language was replaced with baseline-relative evidence,
diagnostic coherence was separated from the selection policy, and the decision
record was defined around target, region, distance, signal source, and quality
flags. These changes improve interpretability without strengthening the
historical outcome evidence. No qualitative feedback archive supports claims
of participant-driven redesign.

The signal-processing and selection rules are deterministic JavaScript
modules. Sixty-three tests cover catalog size and balance, export finiteness
and quality flags, facial-rule behavior, RR filtering, baseline-relative
cardiac behavior, and Random/Vibe selection. They validate the present modules,
not sensor accuracy, Spotify playback, the end-to-end study flow, or the code
run during the pilot.

\subsection{Limitations and threats to validity}

Three limitations define the evidence boundary. First, the facial and cardiac
mappings are fixed heuristics rather than independently validated measures.
Movement may alter camera evidence, raise heart rate, and degrade RR quality at
the same time, allowing one behavior to influence both axes. Short HRV windows
and unmeasured respiration further constrain interpretation.

Second, the small exploratory crossover and incomplete archive prevent
trial-level, order, carryover, missing-data, and subgroup analyses. The
retained files contain no recruitment, ethics, hardware, or trial-level
records, so those parts of the study cannot be reconstructed. One-sided tests
should not be treated as confirmatory without preregistration. The first
adaptive choice may precede completed references, sensing is optional, and the
historical summaries do not identify the signal path behind each choice.

Third, the frozen catalog improves experimental control but does not represent
a dynamic streaming service. The conditions share a catalog, not identical
tracks. The system also uses event-weighted rather than time-weighted
aggregation. At the exact no-signal center, current code paths disagree over
whether the boundary belongs to Calm or Energetic. That edge case should be
unified before another study.

\subsection{Next evaluation}

The next evaluation should begin with controlled changes in expression,
lighting, posture, movement, and respiration applied to the frozen mappings.
An instrumented pilot should then require completed references, timing and
missing-data flags, and a complete ten-trial record. Camera-only,
cardiac-only, fused, and Random conditions would identify which evidence path
contributes. Time-weighted and event-weighted trajectory summaries should be
compared directly.

Before a larger crossover, preregister the smallest relevant mood-fit effect,
sample-size justification, exclusions, primary paired analysis, order
analysis, and signal-quality rules~\cite{lakens2022}. A subsequent
dynamic-catalog study should record candidate availability, playback failures,
coverage, catalog revisions, explanations, override behavior, and privacy
expectations. That sequence separates measurement failure from recommendation
failure and tests whether session-aware re-ranking adds value above a
conventional preference model.
